\section{An anti-invariant module obstruction}
\label{sec:r8-anti-module}
We give eight further exclusions that require no pivotal structure.

\begin{lemma}\label{lem:r8-anti-module}
Let $g$ be an invertible simple of order two in a fusion category, and
let $x_1,\ldots,x_f$ be all simples fixed by right tensoring by $g$,
where $f>0$. There is an integral $f$-dimensional representation $T$ of
the fusion ring such that, with $B_X=(N_{Xx_a}^{x_b})_{a,b}$,
\[
 |(T_X)_{ab}|\leq (B_X)_{ab},\qquad
 (T_X)_{ab}\equiv(B_X)_{ab}\pmod2,\qquad T_{X^*}=T_X^t.
\]
In particular $T_X$ is normal if $XX^*=X^*X$.
\end{lemma}
\begin{proof}
A fixed simple supplies a rank-one module category over the pointed
category generated by $g$, trivializing its associator. Hence the
regular connected separable algebra $A=\mathbf1+g$ exists; compare
\cite[Lemmas~6.1--6.2]{EGODimPQ}. Each fixed simple has two simple right
$A$-module structures $m_a^+,m_a^-$, whereas each free right orbit gives
one simple module. Twisting the right action by the algebra automorphism
that acts as $-1$ on $g$ exchanges $m_a^+$ with $m_a^-$ and fixes the
free-orbit modules. It commutes with left tensoring by every object.
Thus the anti-invariant subgroup has basis
$v_a=[m_a^+]-[m_a^-]$. The two lift multiplicities have sum
$(B_X)_{ab}$ and difference $(T_X)_{ab}$, proving the bound and parity.
Rigidity gives the transpose relation, since the orthogonal difference
vectors all have squared norm two. All fusion relations are preserved.
\end{proof}

\begin{proposition}\label{prop:r8-anti-eight}
\emergencystretch=1em
The rings $R8\_2\_318$, $R8\_2\_319$, $R8\_2\_320$,
$R8\_2\_368$, $R8\_2\_370$, $R8\_2\_735$,
$R8\_2\_758$ and $R8\_2\_760$ have no categorification over $\mathbb C$.
\end{proposition}
\begin{proof}
Set $q=z^2-6z+1$, $r=z^2+2z-1$, and $h=z^2-6z-1$, and put
\[
 A=\begin{pmatrix}1&0&1&1\\0&1&1&1\\1&1&0&1\\1&1&1&2\end{pmatrix},\quad
 B=\begin{pmatrix}0&1&1&1\\1&1&0&1\\1&0&1&1\\1&1&1&2\end{pmatrix},\quad
 C=\begin{pmatrix}0&1&1&1\\1&0&1&1\\1&1&2&1\\1&1&1&2\end{pmatrix}.
\]
The following source identities use zero-based labels. Each indicated
$X$ commutes with $X^*$; the last column is a polynomial vanishing on
$X$. The stated fixed labels are all the right-fixed simples of $g$.
\[
\begin{array}{c|c|c|c|c|c}
\text{suffix}&g&\text{fixed labels}&X&B_X&\text{polynomial}\\\hline
318&4&(1,2,3,7)&5&A&(z^2-1)qr\\
319&4&(1,2,3,7)&5&A&(z^2-1)q\\
320&6&(1,2,5,7)&3&A&(z^2-1)(z^2+1)q\\
368&3&(4,5,6,7)&1&C&(z^2-1)(z^2+1)h\\
370&3&(4,5,6,7)&1&B&(z^2-1)(z^2+1)q\\
735&1&(4,5,6,7)&2&B&(z^2-1)qr\\
758&1&(4,5,6,7)&2&C&(z^2-1)h\\
760&1&(4,5,6,7)&2&B&(z^2-1)q
\end{array}
\]
Let $T=T_X$ from Lemma~\ref{lem:r8-anti-module}. Normality gives
$\|T\|_F^2=\sum_{j=1}^4|\lambda_j|^2$. Integrality forces conjugate
quadratic eigenvalues to occur with equal multiplicity.
For bounds $A,B$, parity gives $\|T\|_F^2\in\{12,16\}$.
The roots $3\pm2\sqrt2$ of $q$ contribute $34$ and so are absent.
For suffixes $319,320,370,760$, all remaining roots have modulus one,
forcing $\|T\|_F^2=4$, a contradiction.

For $318,735$, the additional roots $-1\pm\sqrt2$ contribute $6$.
The possible squared norms are therefore $4,8,12$, according to the
number of copies of this pair. Only two copies are possible, giving
norm squared $12$ and trace $-4$. But norm squared $12$ forces the last
diagonal entry to vanish, while the other diagonal entries are two
signs and a zero. Their trace lies in $[-2,2]$, a contradiction.

For bound $C$, the squared norm is in $\{12,16,20\}$. The roots
$3\pm\sqrt{10}$ of $h$ contribute $38$ and are absent. All remaining
roots for $368,758$ have modulus one, again giving the impossible
squared norm $4$. These arguments give unrestricted exclusions for $318,735$,
in addition to their pivotal obstructions.
\end{proof}

The portable verifier
\texttt{evidence/theoretical\_20261016/rank8/verify\_anti\_module\_eight.py}
checks the displayed source identities and spectral arithmetic directly;
the proof does not rely on an enumeration of candidate matrices.
